La elección de Vue 3 como \textit{framework} de \textit{frontend} responde a su excelente equilibrio entre curva de aprendizaje y potencia. Su sistema de reactividad resulta especialmente adecuado para una aplicación de chat donde los mensajes del tutor aparecen progresivamente en tiempo real, ya que cualquier cambio en los datos se refleja automáticamente en la interfaz. La arquitectura basada en componentes permite organizar la aplicación en componentes independientes, y todo el proyecto utiliza la \textit{Composition API}, la forma moderna de estructurar la lógica en Vue 3. Como empaquetador y servidor de desarrollo se emplea Vite, que destaca por su velocidad de arranque y su sistema de recarga en caliente (\textit{HMR}).

Para el renderizado de las respuestas del tutor, que llegan en formato Markdown con bloques de código, se integran \texttt{marked} y \texttt{highlight.js}. El editor de código se construye sobre Monaco Editor, el mismo motor de Visual Studio Code. La gestión del estado se centraliza con Pinia, los estilos se implementan mediante Tailwind CSS y la comunicación HTTP con el \textit{backend} se realiza a través de Axios.